% GENERATED FILE. Edit canonical lessons, then run npm run generate:books.
% canonical-source-hash: fnv1a64:302fb2d4fe8af10f
% canonical-lessons: FR-C222-societe, FR-C222-gouvernement, FR-C222-loi, FR-C222-tribunal, FR-C222-juge, FR-R222-first-pass-words-from-chapters-214-217, FR-R222-second-pass-words-from-chapters-218-222

\chapter{Society, Government, Law, Court, Judge}
\label{ch:fr-society-government-law-court-judge}
\hlchaptermodality{222}

\begin{chapteropening}
\textbf{By the end of this chapter:} \emph{I can say society, the government, a law, a court and a judge.}

Complete the last lesson of chapter 222: I can say society, the government, a law, a court and a judge.
\end{chapteropening}

\section[la société]{la société — society}

\label{lesson:FR-C222-societe}



\begin{quote}
Before the new one: say the French for a button, then the French for a peach.
\end{quote}

\subsection*{You'll want to know: la société}
\textbf{la société} — \textquotedblleft{}society\textquotedblright{}.

\begin{grammarlens}[title={What you've built}]
One more word for this chapter.
\end{grammarlens}

\subsection*{Your turn}
\begin{itemize}
\raggedright
  \item \emph{Say it:} \emph{la société}
  \item \emph{Say it:} \emph{la société}, once more
  \item \emph{Say it:} say \emph{la société}
  \item \emph{Recall:} say the French for a button, then the French for a peach, then say \emph{la société} again
\end{itemize}

\begin{tcolorbox}[breakable,colback=teal!4,colframe=teal!35!black,title={Before you move on}]
What does la société mean? (\textquotedblleft{}Society\textquotedblright{}.) Say it once more.
\end{tcolorbox}
\section[le gouvernement]{le gouvernement — the government}

\label{lesson:FR-C222-gouvernement}



\begin{quote}
Before the new one: say the French for a peach, then the French for society.
\end{quote}

\subsection*{You'll want to know: le gouvernement}
\textbf{le gouvernement} — \textquotedblleft{}the government\textquotedblright{}.

\begin{grammarlens}[title={What you've built}]
One more word for this chapter.
\end{grammarlens}

\subsection*{Your turn}
\begin{itemize}
\raggedright
  \item \emph{Say it:} \emph{le gouvernement}
  \item \emph{Say it:} \emph{le gouvernement}, once more
  \item \emph{Say it:} say \emph{le gouvernement}
  \item \emph{Recall:} say the French for a peach, then the French for society, then say \emph{le gouvernement} again
\end{itemize}

\begin{tcolorbox}[breakable,colback=teal!4,colframe=teal!35!black,title={Before you move on}]
What does le gouvernement mean? (\textquotedblleft{}The government\textquotedblright{}.) Say it once more.
\end{tcolorbox}
\section[la loi]{la loi — a law}

\label{lesson:FR-C222-loi}



\begin{quote}
Before the new one: say the French for society, then the French for the government.
\end{quote}

\subsection*{You'll want to know: la loi}
\textbf{la loi} — \textquotedblleft{}a law\textquotedblright{}.

\begin{grammarlens}[title={What you've built}]
One more word for this chapter.
\end{grammarlens}

\subsection*{Your turn}
\begin{itemize}
\raggedright
  \item \emph{Say it:} \emph{la loi}
  \item \emph{Say it:} \emph{la loi}, once more
  \item \emph{Say it:} say \emph{la loi}
  \item \emph{Recall:} say the French for society, then the French for the government, then say \emph{la loi} again
\end{itemize}

\begin{tcolorbox}[breakable,colback=teal!4,colframe=teal!35!black,title={Before you move on}]
What does la loi mean? (\textquotedblleft{}A law\textquotedblright{}.) Say it once more.
\end{tcolorbox}
\section[le tribunal]{le tribunal — a court}

\label{lesson:FR-C222-tribunal}



\begin{quote}
Before the new one: say the French for the government, then the French for a law.
\end{quote}

\subsection*{You'll want to know: le tribunal}
\textbf{le tribunal} — \textquotedblleft{}a court\textquotedblright{}.

\begin{grammarlens}[title={What you've built}]
One more word for this chapter.
\end{grammarlens}

\subsection*{Your turn}
\begin{itemize}
\raggedright
  \item \emph{Say it:} \emph{le tribunal}
  \item \emph{Say it:} \emph{le tribunal}, once more
  \item \emph{Say it:} say \emph{le tribunal}
  \item \emph{Recall:} say the French for the government, then the French for a law, then say \emph{le tribunal} again
\end{itemize}

\begin{tcolorbox}[breakable,colback=teal!4,colframe=teal!35!black,title={Before you move on}]
What does le tribunal mean? (\textquotedblleft{}A court\textquotedblright{}.) Say it once more.
\end{tcolorbox}
\section[le juge]{le juge — a judge}

\label{lesson:FR-C222-juge}



\begin{quote}
Before the new one: say the French for a law, then the French for a court.
\end{quote}

\subsection*{You'll want to know: le juge}
\textbf{le juge} — \textquotedblleft{}a judge\textquotedblright{}.

\begin{grammarlens}[title={What you've built}]
One more word for this chapter.
\end{grammarlens}

\subsection*{Your turn}
\begin{itemize}
\raggedright
  \item \emph{Say it:} \emph{le juge}
  \item \emph{Say it:} \emph{le juge}, once more
  \item \emph{Say it:} say \emph{le juge}
  \item \emph{Recall:} say the French for a law, then the French for a court, then say \emph{le juge} again
\end{itemize}

\begin{tcolorbox}[breakable,colback=teal!4,colframe=teal!35!black,title={Before you move on}]
What does le juge mean? (\textquotedblleft{}A judge\textquotedblright{}.) Say it once more.
\end{tcolorbox}
\section[(dialogue)]{Words from chapters 214-217 — a first pass}

\label{lesson:FR-R222-first-pass-words-from-chapters-214-217}



\begin{quote}
Say the words for a court and a judge.
\end{quote}

\subsection*{Your turn}
Say these aloud:

\begin{itemize}
\raggedright
  \item to record, to maintain, to exist, to express, to put on weight
  \item flu, a prescription, a tablet, a syrup, an ointment
  \item around, everywhere, elsewhere, up there, the suburbs
  \item almost, about, only, together, nevertheless
\end{itemize}

\begin{tcolorbox}[breakable,colback=teal!4,colframe=teal!35!black,title={Before you move on}]
To record? (\textbf{enregistrer}.) To maintain? (\textbf{entretenir}.) To exist? (\textbf{exister}.) To express? (\textbf{exprimer}.) To put on weight? (\textbf{grossir}.) Flu? (\textbf{la grippe}.) A prescription? (\textbf{l'ordonnance}.) A tablet? (\textbf{le comprimé}.) A syrup? (\textbf{le sirop}.) An ointment? (\textbf{la pommade}.) Around? (\textbf{autour}.) Everywhere? (\textbf{partout}.) Elsewhere? (\textbf{ailleurs}.) Up there? (\textbf{là-haut}.) The suburbs? (\textbf{la banlieue}.) Almost? (\textbf{presque}.) About? (\textbf{environ}.) Only? (\textbf{seulement}.) Together? (\textbf{ensemble}.) Nevertheless? (\textbf{cependant}.)
\end{tcolorbox}
\section[(dialogue)]{Words from chapters 218-222 — a second pass}

\label{lesson:FR-R222-second-pass-words-from-chapters-218-222}



\begin{quote}
Say the words for curious and jealous.
\end{quote}

\subsection*{Your turn}
Say these aloud:

\begin{itemize}
\raggedright
  \item curious, jealous, proud, nervous, quiet
  \item boring, fun, strange, different, the same
  \item punctual, foreign, friendly, honest, patient
  \item scissors, a watch, pyjamas, a button, a peach
  \item society, the government, a law, a court, a judge
\end{itemize}

\begin{tcolorbox}[breakable,colback=teal!4,colframe=teal!35!black,title={Before you move on}]
Curious? (\textbf{curieux / curieuse}.) Jealous? (\textbf{jaloux / jalouse}.) Proud? (\textbf{fier / fière}.) Nervous? (\textbf{nerveux / nerveuse}.) Quiet? (\textbf{tranquille}.) Boring? (\textbf{ennuyeux / ennuyeuse}.) Fun? (\textbf{amusant / amusante}.) Strange? (\textbf{étrange}.) Different? (\textbf{différent / différente}.) The same? (\textbf{pareil / pareille}.) Punctual? (\textbf{ponctuel / ponctuelle}.) Foreign? (\textbf{étranger / étrangère}.) Friendly? (\textbf{aimable}.) Honest? (\textbf{honnête}.) Patient? (\textbf{patient / patiente}.) Scissors? (\textbf{les ciseaux}.) A watch? (\textbf{la montre}.) Pyjamas? (\textbf{le pyjama}.) A button? (\textbf{le bouton}.) A peach? (\textbf{la pêche}.) Society? (\textbf{la société}.) The government? (\textbf{le gouvernement}.) A law? (\textbf{la loi}.) A court? (\textbf{le tribunal}.) A judge? (\textbf{le juge}.)
\end{tcolorbox}
